\tikzset{
  ctrlBlock/.style={
    draw,
    rounded corners=2pt,
    align=center,
    minimum height=8mm,
    minimum width=24mm,
    inner xsep=4pt,
    inner ysep=3pt,
    font=\small
  },
  ctrlSmall/.style={
    draw,
    rounded corners=2pt,
    align=center,
    minimum height=7mm,
    minimum width=20mm,
    inner xsep=3pt,
    inner ysep=2pt,
    font=\footnotesize
  },
  ctrlSum/.style={
    draw,
    circle,
    inner sep=0pt,
    minimum size=5.5mm,
    font=\small
  },
  ctrlGroup/.style={
    draw,
    dashed,
    rounded corners=3pt,
    inner sep=6pt
  },
  ctrlArrow/.style={
    -{Latex[length=2.2mm]},
    line width=0.45pt
  },
  ctrlDashArrow/.style={
    -{Latex[length=2.2mm]},
    dashed,
    line width=0.45pt
  }
}


\begin{figure}[htbp]
  \centering
  \resizebox{0.98\textwidth}{!}{%
  \begin{tikzpicture}
    \node[ctrlBlock] (task) at (0,0) {链式自重组\\轮足机器人};
    \node[ctrlBlock] (model) at (3.0,0) {单体/链式\\建模};
    \node[ctrlBlock] (single) at (6.0,0) {单体稳定\\控制};
    \node[ctrlBlock] (coop) at (9.0,0) {链式协同\\控制};
    \node[ctrlBlock] (exp) at (12.0,0) {实验验证\\与对比};

    \node[ctrlSmall] (kin) at (3.0,-1.35) {运动学约束};
    \node[ctrlSmall] (lqr) at (6.0,-1.35) {LQR 基础控制};
    \node[ctrlSmall] (dmpc) at (9.0,-1.35) {DMPC / ADRC};

    \draw[ctrlArrow] (task) -- (model);
    \draw[ctrlArrow] (model) -- (single);
    \draw[ctrlArrow] (single) -- (coop);
    \draw[ctrlArrow] (coop) -- (exp);
    \draw[ctrlArrow] (model) -- (kin);
    \draw[ctrlArrow] (single) -- (lqr);
    \draw[ctrlArrow] (coop) -- (dmpc);
    \draw[ctrlDashArrow] (exp.south) -- ++(0,-2.35) -| node[pos=0.28, below, font=\footnotesize] {性能评价与参数修正} (model.south);
  \end{tikzpicture}
  }%
  \caption{第 1 章研究内容与控制方案总体关系框图}
  \label{fig:ch1_control_route}
\end{figure}
